\subsection{Generalized Eigenspaces and Nilpotent Structure}

Throughout this subsection, $\mathcal{V}$ denotes a vector space over $\mathbb{F}$. Generalized eigenvectors enlarge ordinary eigendirections by allowing vectors annihilated by powers of shifted operators.

\begin{definition} \label{def:generalized-eigenvector-and-eigenspace} \hlt{Generalized Eigenvectors and Generalized Eigenspaces}\\
Let $T\in\LL(\mathcal{V})$ and $\lambda\in\mathbb{F}$. Then:
\begin{enumerate}[label=(\roman*)]
\setlength{\itemsep}{0pt}
\item A nonzero $\mathbf{v}\in\mathcal{V}$ is a \hlt{generalized eigenvector} of $T$ corresponding to $\lambda$ iff $(\exists j\in\mathbb{Z}_{>0})\,(T-\lambda I_{\mathcal{V}})^j\mathbf{v}=\mathbf{0}$.
\item The \hlt{generalized eigenspace} corresponding to $\lambda$ is $\mathcal{G}_{\lambda}(T)\defeq\{\mathbf{v}\in\mathcal{V}:(\exists j\in\mathbb{Z}_{>0})\,(T-\lambda I_{\mathcal{V}})^j\mathbf{v}=\mathbf{0}\}$.
\end{enumerate}
\end{definition}

Every eigenvector is generalized by taking $j=1$. For a nonzero generalized eigenvector, successive applications of $T-\lambda I_{\mathcal{V}}$ eventually reach the ordinary eigenspace.

\begin{proposition} \label{prop:kernels-of-powers-increase} \hlt{Kernels of Powers Increase}\\
Let $T\in\LL(\mathcal{V})$, with $T^0\defeq I_{\mathcal{V}}$. Then:
\begin{enumerate}[label=(\roman*)]
\setlength{\itemsep}{0pt}
\item $\ker(T^0)=\{\mathbf{0}\}$;
\item $(\forall k\in\mathbb{Z}_{\geq0})\,\ker(T^k)\subseteq\ker(T^{k+1})$.
\end{enumerate}
\begin{proof}
\begin{enumerate}[label=(\roman*)]
\setlength{\itemsep}{0pt}
\item By $T^0=I_{\mathcal{V}}$, $\ker(T^0)=\{\mathbf{0}\}$.
\item Fix $k\in\mathbb{Z}_{\geq0}$ and $\mathbf{v}\in\ker(T^k)$. Then $T^k\mathbf{v}=\mathbf{0}$, so $T^{k+1}\mathbf{v}=T(T^k\mathbf{v})=\mathbf{0}$. Hence $\mathbf{v}\in\ker(T^{k+1})$.
\end{enumerate}
\hlt{Structural Commentary}: The identity power starts the filtration; one additional application preserves annihilation.\\
\hlt{Key Insight}: An operator cannot revive a vector already sent to zero.
\end{proof}
\end{proposition}

\begin{proposition} \label{prop:kernel-power-stabilization-after-equality} \hlt{Equality Forces Stabilization}\\
Let $T\in\LL(\mathcal{V})$ and $m\in\mathbb{Z}_{\geq0}$. If $\ker(T^m)=\ker(T^{m+1})$, then $(\forall k\in\mathbb{Z}_{\geq0})\,\ker(T^{m+k})=\ker(T^m)$.
\begin{proof}
By Proposition \ref{prop:kernels-of-powers-increase}(ii), it suffices to prove $(\forall k\in\mathbb{Z}_{\geq0})\,\ker(T^{m+k+1})\subseteq\ker(T^{m+k})$.\\
Fix $k\in\mathbb{Z}_{\geq0}$ and $\mathbf{v}\in\ker(T^{m+k+1})$. Then $T^{m+1}(T^k\mathbf{v})=T^{m+k+1}\mathbf{v}=\mathbf{0}$, so $T^k\mathbf{v}\in\ker(T^{m+1})=\ker(T^m)$.\\
Thus $T^{m+k}\mathbf{v}=T^m(T^k\mathbf{v})=\mathbf{0}$, so $\mathbf{v}\in\ker(T^{m+k})$. Hence every adjacent inclusion in the kernel tail is an equality.\\
\hlt{Structural Commentary}: The initial equality converts every later kernel inclusion into equality.\\
\hlt{Key Insight}: Shift a later kernel vector into the already-stable pair.
\end{proof}
\end{proposition}

\begin{theorem} \label{thm:kernel-powers-stabilize-by-dimension} \hlt{Kernels Stop Growing by $\dim(\mathcal{V})$}\\
Let $\mathcal{V}$ be finite-dimensional, let $T\in\LL(\mathcal{V})$, and set $n\defeq\dim(\mathcal{V})$. Then $(\forall k\in\mathbb{Z}_{\geq0})\,\ker(T^{n+k})=\ker(T^n)$.
\begin{proof}
Suppose $\ker(T^n)\neq\ker(T^{n+1})$. By Proposition \ref{prop:kernel-power-stabilization-after-equality}, every inclusion in $\{\mathbf{0}\}=\ker(T^0)\subsetneq\ker(T)\subsetneq\cdots\subsetneq\ker(T^{n+1})$ is strict.\\
Thus $n+1\leq\dim\ker(T^{n+1})\leq n$ by Proposition \ref{prop:kernel-and-range-are-subspaces}(i) and Proposition \ref{prop:dimension-bounds-and-basis-criteria}(i)--(ii), a contradiction.\\
Hence $\ker(T^n)=\ker(T^{n+1})$, and the claim follows by Proposition \ref{prop:kernel-power-stabilization-after-equality}.\\
\hlt{Structural Commentary}: A hypothetical late growth step forces too many strict subspace inclusions.\\
\hlt{Key Insight}: An $n$-dimensional space permits at most $n$ strict kernel increases.
\end{proof}
\end{theorem}

Figure \ref{fig:kernel-power-filtration} records the dimension-limited growth and eventual stabilization of the kernel filtration.

\begin{illustrationfigure}[H]
  \centering
  \begin{tikzpicture}[
  font=\small,
  subspace/.style={draw=blue!60, fill=blue!6, rounded corners=2pt, minimum height=0.65cm, align=center},
  arrow/.style={-{Stealth[length=2mm]}, thick, blue!65}
]
  \node[subspace, minimum width=1.2cm] (k0) at (0,0) {$\ker(T^0)$};
  \node[subspace, minimum width=1.5cm] (k1) at (2.1,0) {$\ker(T)$};
  \node[subspace, minimum width=1.7cm] (k2) at (4.4,0) {$\ker(T^2)$};
  \node[subspace, minimum width=1.2cm] (dots) at (6.5,0) {$\cdots$};
  \node[subspace, minimum width=1.7cm] (kn) at (8.4,0) {$\ker(T^n)$};
  \node[subspace, minimum width=1.8cm] (kn1) at (11.5,0) {$\ker(T^{n+1})$};
  \node[subspace, minimum width=1.2cm] (tail) at (13.5,0) {$\cdots$};

  \draw[arrow] (k0) -- (k1);
  \draw[arrow] (k1) -- (k2);
  \draw[arrow] (k2) -- (dots);
  \draw[arrow] (dots) -- (kn);
  \draw[arrow] (kn) -- node[above, black] {\scriptsize stabilizes} (kn1);
  \draw[arrow] (kn1) -- (tail);

  \draw[gray!70, thick] (-0.7,-0.75) -- (8.75,-0.75);
  \draw[gray!70, thick] (-0.7,-0.65) -- (-0.7,-0.85);
  \draw[gray!70, thick] (8.75,-0.65) -- (8.75,-0.85);
  \node[below, black] at (4.1,-0.95) {at most $\dim(\mathcal{V})$ strict growth steps};
\end{tikzpicture}

  \caption{The kernel filtration grows until it stabilizes; after equality appears once, the tail is constant.}
  \label{fig:kernel-power-filtration}
\end{illustrationfigure}

\begin{theorem} \label{thm:kernel-range-stable-direct-sum} \hlt{Stable Kernel-Range Decomposition}\\
Let $\mathcal{V}$ be finite-dimensional, let $T\in\LL(\mathcal{V})$, and set $n\defeq\dim(\mathcal{V})$. Then $\mathcal{V}=\ker(T^n)\oplus\operatorname{range}(T^n)$.
\begin{proof}
Let $\mathbf{v}\in\ker(T^n)\cap\operatorname{range}(T^n)$. Choose $\mathbf{u}\in\mathcal{V}$ with $\mathbf{v}=T^n\mathbf{u}$ by range membership.\\
Then $T^{2n}\mathbf{u}=T^n\mathbf{v}=\mathbf{0}$ by kernel membership. Thus $\mathbf{u}\in\ker(T^{2n})=\ker(T^n)$ by Theorem \ref{thm:kernel-powers-stabilize-by-dimension}, and $\mathbf{v}=T^n\mathbf{u}=\mathbf{0}$.\\
Hence $\ker(T^n)\cap\operatorname{range}(T^n)=\{\mathbf{0}\}$. By Theorem \ref{thm:rank-nullity}(ii), $\dim\ker(T^n)+\dim\operatorname{range}(T^n)=n$. The sum is direct by Proposition \ref{prop:uniqueness-direct-sum-two-subspaces}.\\
By Corollary \ref{cor:dimension-direct-sum}, $\dim(\ker(T^n)\oplus\operatorname{range}(T^n))=n$. Hence the direct sum equals $\mathcal{V}$ by Proposition \ref{prop:dimension-bounds-and-basis-criteria}(ii).\\
\hlt{Structural Commentary}: Kernel stabilization first kills the intersection; rank--nullity then fills the ambient dimension.\\
\hlt{Key Insight}: Stable image vectors cannot also lie in the stable kernel.
\end{proof}
\end{theorem}

\begin{theorem} \label{thm:generalized-eigenspace-stabilized-kernel} \hlt{Description of Generalized Eigenspaces}\\
Let $\mathcal{V}$ be finite-dimensional, let $T\in\LL(\mathcal{V})$, let $\lambda\in\mathbb{F}$, and set $n\defeq\dim(\mathcal{V})$. Then:
\begin{enumerate}[label=(\roman*)]
\setlength{\itemsep}{0pt}
\item $\mathcal{G}_{\lambda}(T)=\ker((T-\lambda I_{\mathcal{V}})^n)$;
\item $\mathcal{G}_{\lambda}(T)$ is a subspace of $\mathcal{V}$.
\end{enumerate}
\begin{proof}
Set $S\defeq T-\lambda I_{\mathcal{V}}$.
\begin{enumerate}[label=(\roman*)]
\setlength{\itemsep}{0pt}
\item Let $\mathbf{v}\in\mathcal{G}_{\lambda}(T)$. By Definition \ref{def:generalized-eigenvector-and-eigenspace}(ii), choose $j\in\mathbb{Z}_{>0}$ such that $S^j\mathbf{v}=\mathbf{0}$.\\
If $j\leq n$, then $\mathbf{v}\in\ker(S^n)$ by Proposition \ref{prop:kernels-of-powers-increase}(ii); if $j>n$, then $\mathbf{v}\in\ker(S^j)=\ker(S^n)$ by Theorem \ref{thm:kernel-powers-stabilize-by-dimension}. Thus $\mathcal{G}_{\lambda}(T)\subseteq\ker(S^n)$.\\
Conversely, $S^{n+1}\mathbf{v}=\mathbf{0}$ for every $\mathbf{v}\in\ker(S^n)$, so $\ker(S^n)\subseteq\mathcal{G}_{\lambda}(T)$ by Definition \ref{def:generalized-eigenvector-and-eigenspace}(ii).
\item By (i) and Proposition \ref{prop:kernel-and-range-are-subspaces}(i), $\mathcal{G}_{\lambda}(T)=\ker(S^n)$ is a subspace of $\mathcal{V}$.
\end{enumerate}
\hlt{Structural Commentary}: Kernel-chain stabilization replaces all admissible exponents by the single exponent $n$.\\
\hlt{Key Insight}: Every generalized eigenvector is detected by the stabilized shifted kernel.
\end{proof}
\end{theorem}

\begin{proposition} \label{prop:generalized-eigenspaces-invariant} \hlt{Generalized Eigenspaces Are Invariant}\\
Let $\mathcal{V}$ be finite-dimensional and let $T\in\LL(\mathcal{V})$. Then $(\forall\lambda\in\mathbb{F})\,T(\mathcal{G}_{\lambda}(T))\subseteq\mathcal{G}_{\lambda}(T)$.
\begin{proof}
Fix $\lambda\in\mathbb{F}$, set $n\defeq\dim(\mathcal{V})$, and let $p(z)\defeq(z-\lambda)^n$. By Theorem \ref{thm:generalized-eigenspace-stabilized-kernel}(i) and Proposition \ref{prop:operator-polynomial-products}(i), $\mathcal{G}_{\lambda}(T)=\ker(p(T))$.\\
Hence $\mathcal{G}_{\lambda}(T)$ is invariant under $T$ by Proposition \ref{prop:kernel-range-polynomial-in-operator-invariant}(i).\\
\hlt{Structural Commentary}: The stabilized shifted kernel is rewritten as a polynomial kernel in $T$.\\
\hlt{Key Insight}: Polynomial kernels remain invariant under their defining operator.
\end{proof}
\end{proposition}

\begin{definition} \label{def:nilpotent-operator} \hlt{Nilpotent Operator}\\
An operator $N\in\LL(\mathcal{V})$ is \hlt{nilpotent} iff $(\exists m\in\mathbb{Z}_{>0})\,N^m=0$.
\end{definition}

\begin{corollary} \label{cor:nilpotent-power-dimension-zero} \hlt{Nilpotent Power Bound}\\
Let $\mathcal{V}$ be finite-dimensional, let $N\in\LL(\mathcal{V})$ be nilpotent, and set $n\defeq\dim(\mathcal{V})$. Then $N^n=0$.
\begin{proof}
By nilpotence, choose $m\in\mathbb{Z}_{>0}$ such that $N^m=0$, so $\ker(N^m)=\mathcal{V}$. If $m\leq n$, then $\ker(N^m)\subseteq\ker(N^n)$ by Proposition \ref{prop:kernels-of-powers-increase}(ii).\\
If $m>n$, then $\ker(N^m)=\ker(N^n)$ by Theorem \ref{thm:kernel-powers-stabilize-by-dimension}. Thus $\ker(N^n)=\mathcal{V}$ in both cases, and therefore $N^n=0$.\\
\hlt{Structural Commentary}: Kernel growth or stabilization transfers a full nilpotent kernel to exponent $n$.\\
\hlt{Key Insight}: The dimension bound never requires a larger nilpotence exponent.
\end{proof}
\end{corollary}

\begin{lemma} \label{lem:scalar-plus-nilpotent-invertible} \hlt{Nonzero Scalar Plus Nilpotent Is Invertible}\\
Let $N\in\LL(\mathcal{V})$ be nilpotent and let $a\in\mathbb{F}$ with $a\neq0$. Then $aI_{\mathcal{V}}+N$ is invertible.
\begin{proof}
Choose $m\in\mathbb{Z}_{>0}$ such that $N^m=0$ by nilpotence, and set $X\defeq-a^{-1}N$ and $B\defeq a^{-1}\sum\limits_{r=0}^{m-1}X^r$.\\
By Proposition \ref{prop:operator-polynomial-products}(i)--(ii) and the finite geometric-series identity, $(aI_{\mathcal{V}}+N)B=B(aI_{\mathcal{V}}+N)=I_{\mathcal{V}}-X^m=I_{\mathcal{V}}$. Hence $B=(aI_{\mathcal{V}}+N)^{-1}$.\\
\hlt{Structural Commentary}: The finite geometric series verifies both inverse identities for the nilpotent perturbation.\\
\hlt{Key Insight}: Nilpotence makes the geometric-series remainder vanish exactly.
\end{proof}
\end{lemma}

\begin{theorem} \label{thm:distinct-generalized-eigenvectors-independent} \hlt{Generalized Eigenvectors for Distinct Eigenvalues Are Independent}\\
Let $\mathcal{V}$ be finite-dimensional, let $T\in\LL(\mathcal{V})$, and let $m\in\mathbb{Z}_{>0}$. Suppose $\lambda_1,\ldots,\lambda_m$ are distinct eigenvalues of $T$, and let $\mathbf{v}_1,\ldots,\mathbf{v}_m$ be corresponding generalized eigenvectors.\\
Then $(\mathbf{v}_1,\ldots,\mathbf{v}_m)$ is linearly independent.
\begin{proof}
\begin{enumerate}[label=(\roman*)]
\setlength{\itemsep}{0pt}
\item \hlt{Candidate Relation}: Set $n\defeq\dim(\mathcal{V})$ and suppose $a_1\mathbf{v}_1+\cdots+a_m\mathbf{v}_m=\mathbf{0}$.
\item \hlt{Polynomial Filter}: Fix $r\in\{1,\ldots,m\}$ and set $Q_r(T)\defeq\prod\limits_{i\in\{1,\ldots,m\}\setminus\{r\}}(T-\lambda_iI_{\mathcal{V}})^n$, with the empty product equal to $I_{\mathcal{V}}$.
\item \hlt{Non-Target Annihilation}: By Theorem \ref{thm:generalized-eigenspace-stabilized-kernel}(i) and commutativity in Proposition \ref{prop:operator-polynomial-products}(ii), $Q_r(T)\mathbf{v}_i=\mathbf{0}$ for $i\neq r$.\\
Hence applying $Q_r(T)$ to the candidate relation leaves $a_rQ_r(T)\mathbf{v}_r=\mathbf{0}$.
\item \hlt{Target Invertibility}: Set $\mathcal{G}_r\defeq\mathcal{G}_{\lambda_r}(T)$ and $N_r\defeq(T-\lambda_rI_{\mathcal{V}})|_{\mathcal{G}_r}$.\\
By Proposition \ref{prop:generalized-eigenspaces-invariant} and Theorem \ref{thm:generalized-eigenspace-stabilized-kernel}(i)--(ii), $N_r\in\LL(\mathcal{G}_r)$ and $N_r^n=0$.\\
For $i\neq r$, $(T-\lambda_iI_{\mathcal{V}})|_{\mathcal{G}_r}=N_r+(\lambda_r-\lambda_i)I_{\mathcal{G}_r}$ with $\lambda_r-\lambda_i\neq0$, so this restriction is invertible by Lemma \ref{lem:scalar-plus-nilpotent-invertible}.\\
Thus $Q_r(T)|_{\mathcal{G}_r}$, a product of $n$th powers of these restrictions or the identity if $m=1$, is invertible by Proposition \ref{prop:inverse-of-product-linear-maps}(i).
\item \hlt{Coefficient Vanishing}: Since $\mathbf{v}_r\in\mathcal{G}_r\setminus\{\mathbf{0}\}$, $Q_r(T)\mathbf{v}_r\neq\mathbf{0}$ by item (iv), and hence $a_r=0$ by item (iii).\\
Since $r$ was arbitrary, $a_1=\cdots=a_m=0$, so $(\mathbf{v}_1,\ldots,\mathbf{v}_m)$ is linearly independent.
\end{enumerate}
\hlt{Structural Commentary}: Polynomial filters annihilate non-target vectors while nilpotent-plus-scalar factors preserve the target.\\
\hlt{Key Insight}: Distinct spectral shifts are invertible on the remaining generalized eigenspace.
\end{proof}
\end{theorem}

\begin{theorem} \label{thm:nilpotent-strictly-upper-triangular-matrix} \hlt{Nilpotent Operators and Strictly Upper-Triangular Matrices}\\
Let $\mathcal{V}$ be nonzero and finite-dimensional, let $N\in\LL(\mathcal{V})$, and set $n\defeq\dim(\mathcal{V})$.\\
Then $N$ is nilpotent iff some ordered basis $\beta$ satisfies $[N]_{\beta}=[a_{ij}]$ and $(\forall i,j\in\{1,\ldots,n\})\,(i\geq j\implies a_{ij}=0)$.
\begin{proof}
\begin{enumerate}[label=\roman*.]
\setlength{\itemsep}{0pt}
\item[$(\Rightarrow)$] Suppose $N$ is nilpotent. By Corollary \ref{cor:nilpotent-power-dimension-zero} and Proposition \ref{prop:kernels-of-powers-increase}(i)--(ii), $\{\mathbf{0}\}=\ker(N^0)\subseteq\ker(N)\subseteq\cdots\subseteq\ker(N^n)=\mathcal{V}$.\\
By Corollary \ref{cor:finite-dimensional-spaces-and-subspaces-have-bases}(i)--(ii) and Theorem \ref{thm:extending-independent-lists}, choose an ordered basis $\beta$ whose initial segments span the kernels in this chain.\\
If $\mathbf{v}$ first enters at level $k$, then $N\mathbf{v}\in\ker(N^{k-1})$, which is spanned by earlier vectors of $\beta$; hence $[N]_{\beta}$ is strictly upper triangular.
\item[$(\Leftarrow)$] Suppose $\beta=(\mathbf{v}_1,\ldots,\mathbf{v}_n)$ and $[N]_{\beta}$ is strictly upper triangular.\\
Set $F_0\defeq\{\mathbf{0}\}$ and $F_j\defeq\operatorname{span}(\mathbf{v}_1,\ldots,\mathbf{v}_j)$ for $j\in\{1,\ldots,n\}$. Then $(\forall j\in\{1,\ldots,n\})\,N(F_j)\subseteq F_{j-1}$.\\
By iteration, $N^n(\mathcal{V})=N^n(F_n)\subseteq F_0=\{\mathbf{0}\}$, so $N^n=0$ and $N$ is nilpotent by Definition \ref{def:nilpotent-operator}.
\end{enumerate}
\hlt{Structural Commentary}: The forward direction adapts a basis to kernels; the reverse direction descends basis indices.\\
\hlt{Key Insight}: Strict index descent forces every basis vector to reach zero.
\end{proof}
\end{theorem}

Figure \ref{fig:nilpotent-filtration-basis} illustrates how a basis adapted to the kernel filtration makes nilpotent action move toward earlier layers.

\begin{illustrationfigure}[H]
  \centering
  \begin{tikzpicture}[
  font=\small,
  layer/.style={draw=teal!65, fill=teal!7, rounded corners=2pt, minimum width=1.95cm, minimum height=0.55cm, align=center},
  arrow/.style={-{Stealth[length=2mm]}, thick, teal!70}
]
  \node[layer] (n1) at (0,0) {$\ker(N)$};
  \node[layer] (n2) at (2.46,0) {$\ker(N^2)$};
  \node[layer] (n3) at (4.93,0) {$\ker(N^3)$};
  \node[layer] (dots) at (7.13,0) {$\cdots$};
  \node[layer] (nn) at (9.42,0) {$\ker(N^n)=\mathcal{V}$};

  \draw[arrow] (n1) -- (n2);
  \draw[arrow] (n2) -- (n3);
  \draw[arrow] (n3) -- (dots);
  \draw[arrow] (dots) -- (nn);
\end{tikzpicture}

  \caption{Extending bases through the kernel filtration makes $N$ move vectors to earlier layers.}
  \label{fig:nilpotent-filtration-basis}
\end{illustrationfigure}

\begin{perspective} \label{per:generalized-eigenvectors-as-nilpotent-data} \hlt{Eigenvalue Plus Nilpotent Part}\\
Let $\mathcal{V}$ be finite-dimensional, let $T\in\LL(\mathcal{V})$, let $\lambda\in\mathbb{F}$, and set $n\defeq\dim(\mathcal{V})$. Define $\mathcal{G}_{\lambda}\defeq\mathcal{G}_{\lambda}(T)$ and $N_{\lambda}\defeq(T-\lambda I_{\mathcal{V}})|_{\mathcal{G}_{\lambda}}$.\\
By Proposition \ref{prop:generalized-eigenspaces-invariant}, $N_{\lambda}\in\LL(\mathcal{G}_{\lambda})$. By Theorem \ref{thm:generalized-eigenspace-stabilized-kernel}(i), $N_{\lambda}^n=0$, and hence $N_{\lambda}^{n+1}=0$. Thus $N_{\lambda}$ is nilpotent and $T|_{\mathcal{G}_{\lambda}}=\lambda I_{\mathcal{G}_{\lambda}}+N_{\lambda}$.\\
Generalized eigenvectors therefore record the nilpotent component of $T$ within the stabilized shifted kernel.
\end{perspective}
